\chapter{Token-Management -- Der heimliche Kostenfaktor}
\label{chap:token_management}

% ============================================================
% LERNZIELE
% ============================================================
\begin{lernziele}
Nach diesem Kapitel kannst du:
\begin{itemize}
    \item Die Tokenisierung für GPT, Claude und Gemini im Detail erklären und deutsche Token-Effizienz berechnen
    \item TikToken und alternative Tokenizer in Python produktiv einsetzen
    \item Prompt-Kompaktierung, Chat-History-Management und Prompt-Caching strategisch kombinieren
    \item Token-Budgets für Produktionssysteme entwerfen und Cost-Forecasts erstellen
    \item Multi-modale Token-Kosten (Bilder, Audio, Video) korrekt kalkulieren
\end{itemize}
\end{lernziele}

% ============================================================
% MOTIVATION
% ============================================================
\section{Motivation}

Token-Management ist die grundlegendste Engineering-Disziplin im Umgang mit LLMs. Während Software-Engineer:innen normalerweise in Request-Volumen oder Datenvolumen denken, wird bei LLMs in \textbf{Tokens abgerechnet}. Ein Token ist die kleinste Verrechnungseinheit -- kein Request, kein Byte, kein Wort.

Die Konsequenz: Ohne Token-Management weißt du weder, was ein API-Call kostet, noch, ob der Prompt ins Context-Limit passt. In der Praxis führen falsche Token-Kalkulationen zu:
\begin{itemize}[leftmargin=*]
    \item \textbf{Kosten-Explosion}: Ein Batch-Job, der auf 10K Requests ausgelegt war, kostet plötzlich das Zehnfache, weil jeder Prompt 4.000 statt 400 Tokens hatte.
    \item \textbf{Cutoff-Fehler}: Requests werden stillschweigend abgeschnitten, weil der Prompt das Context-Limit überschreitet -- der Fehler erscheint erst in der Response.
    \item \textbf{Inkonsistente Ergebnisse}: Unterschiedliche Tokenizer (OpenAI vs. Anthropic) liefern unterschiedliche Resultate bei gleichem Text.
\end{itemize}

Token-Management ist keine optionale Optimierung, sondern eine \textbf{Grundvoraussetzung} für jede produktive LLM-Integration.

% ============================================================
% GRUNDLAGEN
% ============================================================
\section{Grundlagen -- Tokenisierung im Detail}

\subsection{Was ist ein Token?}

Ein Token ist die kleinste Einheit, die ein LLM verarbeiten kann. Anders als natürliche Sprache, die in Wörtern oder Zeichen denkt, arbeitet ein Transformer auf einer \textbf{vordefinierten Vokabelliste} von Tokens. Ein Token kann sein:
\begin{itemize}[leftmargin=*]
    \item Ein ganzes Wort (z.\,B. \texttt{the} = 1 Token)
    \item Ein Wortteil (z.\,B. \texttt{un} + \texttt{believe} + \texttt{able})
    \item Ein einzelnes Zeichen bei unbekannten Wörtern
    \item Whitespace und Satzzeichen (jeweils eigene Tokens)
\end{itemize}

Die Umwandlung von Text in Tokens heißt \textbf{Tokenisierung}. Der Tokenizer ist ein fester Bestandteil jedes Modells -- er wird beim Training mitgelernt und ändert sich nicht mehr.

\subsection{Die drei Haupt-Tokenisierer}

\begin{table}[H]
\centering
\begin{tabularx}{\linewidth}{lXX}
\toprule
\textbf{Modell-Familie} & \textbf{Tokenizer} & \textbf{Algorithmus} \\
\midrule
GPT-4o / GPT-4 / o1 & TikToken & BPE (Byte-Pair-Encoding) \\
Claude 3.5 / Claude 4 & Claude Tokenizer & SentencePiece (BPE) \\
Gemini 2.0 / 2.5 & Google SentencePiece & SentencePiece (Unigram) \\
LLaMA 3.x & LLaMA Tokenizer & BPE (gguf-optimiert) \\
\midrule
\end{tabularx}
\caption{Tokenizer der wichtigsten Modell-Familien}
\label{tab:tokenizer-vergleich}
\end{table}

\subsubsection{BPE (Byte-Pair-Encoding)}

BPE ist der dominierende Tokenisierungs-Algorithmus. Er funktioniert in drei Schritten:
\begin{enumerate}[leftmargin=*]
    \item \textbf{Zerlegung}: Jedes Zeichen wird einzeln als Token behandelt
    \item \textbf{Merging}: Die häufigsten Paare werden iterativ zu einem Token zusammengefasst
    \item \textbf{Vokabular}: Nach einer festgelegten Anzahl von Merges entsteht ein Vokabular (GPT-4o: ca. 100K Tokens)
\end{enumerate}

\subsubsection{SentencePiece}

SentencePiece ist ein \textbf{sprachunabhängiger} Tokenizer, der direkt auf rohem Text arbeitet, ohne Whitespace-Pre-Tokenisierung. Er wird von Anthropic und Google verwendet.

\subsection{Tokenisierung im Vergleich: Deutsch vs. Englisch}

Deutsch erzeugt typischerweise \textbf{1,4--2x mehr Tokens als Englisch} für denselben Inhalt. Der Grund liegt in deutschen Komposita:

\begin{table}[H]
\centering
\begin{tabularx}{\linewidth}{lXXX}
\toprule
\textbf{Wort} & \textbf{GPT-4o (BPE)} & \textbf{Tokens} & \textbf{Bemerkung} \\
\midrule
the cat & [the] [cat] & 2 & Kurze Wörter \\
Kraftfahrzeughaftpflicht & [K] [raft] [fahr] [zeug] [haf] [t] [pfl] [icht] & 8 & Kompositum zerfällt \\
Haftpflichtversicherung & [Haf] [t] [pfl] [icht] [vers] [ich] [er] [ung] & 8 & Langes Wort \\
\midrule
\end{tabularx}
\caption{Token-Zerlegung Deutsch vs. Englisch (gpt-4o)}
\label{tab:token-deutsch}
\end{table}

Ein 500-Wort-englischer Text wird typischerweise zu \textbf{\~650 Tokens}, der gleiche deutsche Text zu \textbf{\~850--1.000 Tokens}. Das bedeutet 30--50\,\% höhere Kosten für deutsche Texte bei gleichem Inhalt.

\subsection{Tokenisierung visualisieren}

\begin{verbatim}
Text: "Haftpflichtversicherung ist ein langes Wort."

GPT-4o (tiktoken):
  [Haf] [t] [pfl] [icht] [vers] [ich] [er] [ung] [ist] [ein] [lang] [es] [Wort] [.]
  = 14 Tokens

Claude (SentencePiece):
  [_Haft] [pfl] [icht] [vers] [icher] [ung] [_ist] [_ein] [_langes] [_Wort] [.]
  = 11 Tokens (effizienter bei Komposita)
\end{verbatim}

\begin{praxishinweis}
Deutsch produziert 30--50\,\% mehr Tokens als Englisch für denselben Inhalt. Das ist kein Bug, sondern ein Feature der Tokenizer -- muss aber in der Kostenkalkulation berücksichtigt werden. Besonders RAG-Systeme mit deutschen Texten sind stärker betroffen als reine Chat-Anwendungen.
\end{praxishinweis}

% ============================================================
% ARCHITEKTUR
% ============================================================
\section{Architektur -- Wie Tokens durch die Pipeline fließen}

Token-Management ist kein isolierter Schritt, sondern durchzieht die gesamte LLM-Pipeline:

\begin{verbatim}
  Text -> Tokenizer -> Embedding-Lookup -> Transformer -> Output-Logits -> Detokenizer -> Text

  Kosten-Falle #1:        Kosten-Falle #2:         Kosten-Falle #3:
  Unnötig lange Prompts   Kein max_tokens          Chat-History wächst unbegrenzt
\end{verbatim}

\subsection{Kosten-Arten im Token-Fluss}

\begin{table}[H]
\centering
\begin{tabularx}{\linewidth}{lXl}
\toprule
\textbf{Kostenart} & \textbf{Entstehung} & \textbf{Hebel} \\
\midrule
Input-Tokens & Prompt (System + User + History) & Prompt-Kompaktierung, Caching \\
Output-Tokens & Modell-Antwort & max\_tokens, kurze Antworten erzwingen \\
Overhead-Tokens & API-Formatierung (Rollen, Tool-Defs) & Weniger Tools, kürzere Tool-Beschreibungen \\
Multi-modal & Bilder, Audio, Video im Prompt & Auflösung reduzieren, Audio komprimieren \\
Batch-Processing & Wiederholte Requests mit gleichem Prefix & Prompt-Caching \\
\midrule
\end{tabularx}
\caption{Kostenarten im Token-Fluss}
\label{tab:kostenarten}
\end{table}

\subsection{Token-Overhead pro Provider}

Jeder API-Provider fügt Overhead-Tokens hinzu, die bei der Kostenkalkulation oft vergessen werden:

\begin{table}[H]
\centering
\begin{tabularx}{\linewidth}{lX}
\toprule
\textbf{Provider} & \textbf{Overhead pro Message} \\
\midrule
OpenAI & 3 Tokens pro Message + 1 Token pro Name \\
Anthropic & ~4 Tokens pro Message + System-Prompt-Separator \\
Google & ~3 Tokens pro Message + Content-Blöcke \\
\midrule
\end{tabularx}
\caption{Message-Overhead pro Provider}
\label{tab:overhead-provider}
\end{table}

% ============================================================
% THEORIE
% ============================================================
\section{Theorie -- Warum Token-Management nicht trivial ist}

\subsection{Das Invarianz-Problem}

Token-ization ist \textbf{nicht invariant}. Derselbe Semantik-Inhalt kann in sehr unterschiedlich vielen Tokens kodiert sein:

\begin{itemize}[leftmargin=*]
    \item \textbf{Umformulierung}: "Der Hund läuft schnell" vs. "Der flinke Hund rennt" -- gleiche Semantik, unterschiedliche Token-Zahl
    \item \textbf{Whitespace}: Ein zusätzliches Leerzeichen oder Zeilenumbruch erzeugt einen zusätzlichen Token
    \item \textbf{JSON-Formatierung}: Gleiches JSON formatiert vs. minimiert unterscheidet sich um 20--40\,\% Tokens
\end{itemize}

\subsection{Das nicht-lineare Kosten-Verhältnis}

Die Kosten steigen nicht linear mit der Token-Zahl, weil:
\begin{itemize}[leftmargin=*]
    \item \textbf{Pricing-Stufen}: Input kostet weniger als Output (Faktor 3--10x)
    \item \textbf{Context-Length-Pricing}: Längere Kontexte haben höhere Token-Preise (Cache- vs. Non-Cache-Rate)
    \item \textbf{Batch-Rabatt}: Batch-Processing kostet 50\,\% weniger
\end{itemize}

Ein Prompt mit 100K Input-Tokens und 100 Output-Tokens kostet bei GPT-4o:
\begin{itemize}[leftmargin=*]
    \item Input: 100.000 $\times$ \$2,50 / 1M = \$0,25
    \item Output: 100 $\times$ \$10,00 / 1M = \$0,001
    \item Gesamt: \textbf{\$0,251}
\end{itemize}

Gleiches Verhältnis umgekehrt (100 Input, 100K Output):
\begin{itemize}[leftmargin=*]
    \item Input: 100 $\times$ \$2,50 / 1M = \$0,00025
    \item Output: 100.000 $\times$ \$10,00 / 1M = \$1,00
    \item Gesamt: \textbf{\$1,00}
\end{itemize}

\textbf{Fazit: Output-Tokens dominieren die Kosten bei langen Antworten, Input-Tokens bei RAG-Systemen mit großen Kontexten.}

\begin{praxishinweis}
Die Kostenasymmetrie zwischen Input und Output (Faktor 3--10x) bestimmt die optimale Architektur: Für RAG-Systeme mit vielen Input- und wenigen Output-Tokens lohnt sich Prompt-Caching enorm. Für Chat-Systeme mit langen Antworten ist max\_tokens der wichtigste Hebel.
\end{praxishinweis}

% ============================================================
% PRAXIS
% ============================================================
\section{Praxisbeispiele aus der Industrie}

\subsection{Fintech: Transaktions-Kategorisierung}

Ein Fintech-Unternehmen kategorisiert 500.000 Transaktionsbeschreibungen monatlich:

\begin{itemize}[leftmargin=*]
    \item \textbf{Prompt}: 50 Tokens (System + Transaktions-Text)
    \item \textbf{Output}: 10 Tokens (Kategorie + Confidence)
    \item \textbf{Gesamt pro Request}: 60 Tokens
    \item \textbf{Monatsvolumen}: 500.000 $\times$ 60 = 30M Tokens
    \item \textbf{Kosten GPT-4o-mini}: \$0,45 (Input) + \$0,15 (Output) = \textbf{\$0,60}
    \item \textbf{Kosten GPT-4o}: \$4,50 (Input) + \$1,50 (Output) = \textbf{\$6,00}
\end{itemize}

Die Wahl des Modells bestimmt hier die Kosten um den Faktor 10.

\subsection{E-Commerce: Chat-Support mit History-Management}

Ein Online-Shop speichert den gesamten Chat-Verlauf pro Kunde:

\begin{itemize}[leftmargin=*]
    \item \textbf{Ohne Kompression}: Nach 10 Nachrichten $\rightarrow$ 4.500 Tokens $\rightarrow$ Kosten pro Chat: \$0,11 (GPT-4o)
    \item \textbf{Mit Kompression}: Nach 10 Nachrichten $\rightarrow$ 250 Tokens (zusammengefasst) $\rightarrow$ Kosten pro Chat: \$0,01
    \item \textbf{Täglich}: 5.000 Chats $\rightarrow$ \$50 vs. \$550 pro Tag
\end{itemize}

Die History-Kompression spart hier 90\,\% der Kosten.

\subsection{SaaS: Prompt-Caching in der Praxis}

Ein SaaS-Anbieter hat einen 4K-System-Prompt, der bei jedem Call mitgesendet wird:

\begin{itemize}[leftmargin=*]
    \item \textbf{100.000 Requests / Tag}, 4K Input + 500 Output
    \item \textbf{Ohne Caching}: 100K $\times$ (4K + 0,5K) = 450M Tokens $\times$ GPT-4o-Preis
    \item \textbf{Mit Caching}: 100K $\times$ (0K Cache + 0,5K Output) = 50M Tokens
    \item \textbf{Ersparnis}: 89\,\% der Input-Kosten = \textbf{\$400/Tag vs. \$45/Tag}
\end{itemize}

\begin{praxishinweis}
Die drei Praxisbeispiele zeigen ein Muster: In jedem Fall ist es die Kombination aus Modellwahl, History-Management und Caching, die die Kosten drastisch senkt. Keine einzelne Optimierung bringt 90\,\% -- erst das Zusammenspiel aller drei Hebel.
\end{praxishinweis}

% ============================================================
% CODE: TIKTOKEN (existiert, erweitert)
% ============================================================
\section{Token-Counting in Python}

Der wichtigste Schritt zur Kostenkontrolle ist, Prompts vor dem Senden zu zählen. Hier die produktive Implementierung.

\subsection{TikToken -- der OpenAI-Standard}

\label{sec:tiktoken}

\begin{lstlisting}[language=Python, caption=TikToken: Tokens vor dem API-Call zaehlen]
import tiktoken

MODEL_ENCODINGS = {
    "gpt-4o": "o200k_base",
    "gpt-4o-mini": "o200k_base",
    "o1": "o200k_base",
    "o3": "o200k_base",
    "gpt-4": "cl100k_base",
    "gpt-3.5-turbo": "cl100k_base",
}

def get_encoding(model: str) -> tiktoken.Encoding:
    encoding_name = MODEL_ENCODINGS.get(model, "o200k_base")
    return tiktoken.get_encoding(encoding_name)

def count_tokens(messages: list[dict], model: str = "gpt-4o") -> int:
    encoding = get_encoding(model)
    tokens_per_message = 3
    tokens_per_name = 1

    total = 0
    for msg in messages:
        total += tokens_per_message
        for key, value in msg.items():
            if isinstance(value, str):
                total += len(encoding.encode(value))
            elif isinstance(value, list):
                for item in value:
                    if item.get("type") == "text":
                        total += len(encoding.encode(item["text"]))
                    elif item.get("type") == "image_url":
                        total += count_image_tokens(item["image_url"]["url"])
            if key == "name":
                total += tokens_per_name

    total += 3
    return total

def count_image_tokens(image_url: str, detail: str = "auto") -> int:
    if detail == "low":
        return 85
    if detail == "high":
        return 170  # 768x768 default
    if detail == "auto":
        return 85
    return 85
\end{lstlisting}

\subsection{Alternative Provider-Tokenizer}

Nicht nur OpenAI benötigt Token-Counting:

\begin{lstlisting}[language=Python, caption=Anthropic Token-Counting (Claude)]
# Anthropic: count_tokens ist im SDK enthalten
from anthropic import Anthropic

client = Anthropic()
token_count = client.count_tokens(
    "Das ist ein Test-Prompt fuer Claude."
)
print(f"Claude Tokens: {token_count}")

# Oder manuell mit anthropic-tokenizer
from anthropic_tokenizer import count_tokens as claude_count

tokens = claude_count("Dein Prompt hier")
\end{lstlisting}

\begin{lstlisting}[language=Python, caption=Google Gemini Token-Counting]
# Google: Token-Counting ueber die API
import google.generativeai as genai

model = genai.GenerativeModel("gemini-2.0-flash")
response = model.count_tokens("Dein Prompt hier")
print(f"Gemini Tokens: {response.total_tokens}")
\end{lstlisting}

\subsection{Multi-Provider Token-Counter}

Für Produktionssysteme mit mehreren Providern:

\begin{lstlisting}[language=Python, caption=Unabhaengiger Token-Counter fuer alle Provider]
class TokenCounter:
    def __init__(self):
        self._openai_encoding = tiktoken.get_encoding("o200k_base")

    def count_openai(self, text: str) -> int:
        return len(self._openai_encoding.encode(text))

    def count_anthropic(self, text: str) -> int:
        from anthropic import Anthropic
        return Anthropic().count_tokens(text)

    def count_gemini(self, text: str, model_name: str = "gemini-2.0-flash") -> int:
        import google.generativeai as genai
        model = genai.GenerativeModel(model_name)
        return model.count_tokens(text).total_tokens

    def count_llama(self, text: str) -> int:
        import tiktoken
        encoding = tiktoken.get_encoding("cl100k_base")
        return len(encoding.encode(text))

    def estimate_cost(self, provider: str, model: str,
                      input_tokens: int, output_tokens: int) -> float:
        prices = {
            ("openai", "gpt-4o"):       (2.50, 10.00),
            ("openai", "gpt-4o-mini"):  (0.15, 0.60),
            ("anthropic", "claude-4"):  (3.00, 15.00),
            ("google", "gemini-2.0"):   (1.00, 4.00),
        }
        input_price, output_price = prices.get(
            (provider, model), (0.0, 0.0)
        )
        return (input_tokens / 1_000_000 * input_price +
                output_tokens / 1_000_000 * output_price)
\end{lstlisting}

\section{Max Tokens verstehen und konfigurieren}

Der \texttt{max\_tokens}-Parameter kontrolliert nicht nur die Antwortlänge, sondern ist auch ein \textbf{kritischer Kostenfaktor}:

\begin{itemize}[leftmargin=*]
    \item \textbf{Ohne max\_tokens:} Das Modell kann bis zum Context-Limit generieren -- bei komplexen Queries schnell 4096 Tokens (und damit €0,06+ nur für den Output bei Claude Sonnet).
    \item \textbf{Mit max\_tokens:} Setze einen vernünftigen Wert basierend auf der erwarteten Antwortlänge. Eine Klassifizierung braucht 50 Tokens, eine Analyse 500.
    \item \textbf{Output-Overhead:} Plane immer 10\,\% Puffer für den internen Output-Zähler ein.
\end{itemize}

\begin{lstlisting}[language=Python, caption={max\_tokens sinnvoll setzen}]
TASK_MAX_TOKENS = {
    "classification":     50,
    "extraction":        200,
    "summarization":     512,
    "analysis":         1024,
    "code_generation":  2048,
    "creative_writing": 4096,
}

def get_max_tokens(task_type: str) -> int:
    return TASK_MAX_TOKENS.get(task_type, 512)

# Dynamisches max_tokens
def configure_request(task_type: str, user_input: str) -> dict:
    input_tokens = count_tokens(user_input)
    max_out = get_max_tokens(task_type)
    return {
        "max_tokens": max_out,
        "estimated_cost": estimate_cost(
            "openai", "gpt-4o", input_tokens, max_out
        ),
    }
\end{lstlisting}

% ============================================================
% BEST PRACTICES
% ============================================================
\section{Best Practices}

\subsection{Prompt-Kompaktierung -- die wichtigste Kostenoptimierung}

Der System-Prompt wird bei \textbf{jeder API-Anfrage mitgesendet}. Ein 4K-Token-System-Prompt kostet also bei 100K Requests: 400M Input-Tokens. Hier die effektivsten Kompaktierungsstrategien:

\begin{lstlisting}[language=Python, caption=Prompt-Kompaktierung]
# Schlecht: Langer, wiederholender Prompt
SYSTEM_PROMPT_LANG = """Du bist ein Kundenservice-Assistent fuer eine
E-Commerce-Firma. Bitte antworte immer hoeflich und professionell.
Deine Aufgabe ist es, Kundenanfragen zu beantworten. Wenn du etwas
nicht weisst, sage das. Vermeide lange Antworten. Sei kurz und
praezise. Antworte auf Deutsch."""

# Gut: Kompaktierter Prompt (70% weniger Tokens)
SYSTEM_PROMPT_KURZ = (
    "Du bist E-Commerce-Kundenservice. "
    "Antworte hoeflich, praezise, auf Deutsch. "
    "Bei Unsicherheit: nachfragen."
)
\end{lstlisting}

\subsection{Chat-History-Management}

Die Chat-History ist der größte versteckte Kostentreiber:

\begin{lstlisting}[language=Python, caption=Chat-History mit automatischer Kompression]
from collections import deque
import tiktoken

class ConversationMemory:
    def __init__(self, model: str = "gpt-4o", max_tokens: int = 120_000):
        self.model = model
        self.encoding = tiktoken.encoding_for_model(model)
        self.max_tokens = max_tokens
        self.history: deque[dict] = deque()
        self.system_prompt: str = ""

    def add_message(self, role: str, content: str):
        if role == "system":
            self.system_prompt = content
            return
        self.history.append({"role": role, "content": content})
        if self._current_tokens() > self.max_tokens * 0.8:
            self._compress()

    def _current_tokens(self) -> int:
        total = len(self.encoding.encode(self.system_prompt))
        for msg in self.history:
            total += 3 + len(self.encoding.encode(msg["content"]))
        return total

    def _compress(self):
        if len(self.history) < 4:
            return
        recent = list(self.history)[-2:]
        old = list(self.history)[:-2]
        summary = openai_client.chat.completions.create(
            model="gpt-4o-mini",
            messages=[{
                "role": "user",
                "content": "Fasse diese Chat-Nachrichten in einem Satz"
                           " zusammen: " + " ".join(
                               m["content"] for m in old
                           )
            }],
            max_tokens=100,
        ).choices[0].message.content
        self.history = deque(
            [{"role": "user", "content": f"[Summary: {summary}]"}]
            + recent
        )

    def get_messages(self) -> list[dict]:
        msgs = [{"role": "system", "content": self.system_prompt}]
        msgs.extend(self.history)
        return msgs
\end{lstlisting}

\subsection{Prompt-Caching}

Prompt-Caching ist der \textbf{wichtigste Kostenhebel}. Sowohl OpenAI als auch Anthropic unterstützen es:

\begin{lstlisting}[language=Python, caption=Prompt-Caching mit OpenAI]
# OpenAI: Automatisches Caching fuer lange System-Prompts
# Caching passiert automatisch ab 1024 Token Input
# Gecachte Tokens kosten ~50% weniger

response = client.chat.completions.create(
    model="gpt-4o",
    messages=[
        {"role": "system", "content": system_prompt},  # Gecached!
        {"role": "user", "content": user_input},
    ],
)

# Nutze response.usage fuer Cache-Statistiken
usage = response.usage
print(f"Cache Read: {usage.prompt_tokens_details.cached_tokens}")
print(f"Tokens Input (ohne Cache): {usage.prompt_tokens}")
\end{lstlisting}

\begin{lstlisting}[language=Python, caption=Prompt-Caching mit Anthropic]
# Anthropic: Explizites Cache-Control (ab Claude 3.5)
response = client.messages.create(
    model="claude-3-5-sonnet-20241022",
    system=[{
        "type": "text",
        "text": system_prompt,
        "cache_control": {"type": "ephemeral"},
    }],
    messages=[{"role": "user", "content": user_input}],
)

# Cache-Statistiken auslesen
print(f"Cache Created: {response.usage.cache_creation_input_tokens}")
print(f"Cache Read: {response.usage.cache_read_input_tokens}")
\end{lstlisting}

\subsection{Token-Budget-System}

Definiere ein Budget pro Request, bevor du die API aufrufst:

\begin{lstlisting}[language=Python, caption=Token-Budget-System]
class TokenBudget:
    def __init__(self, max_input: int, max_output: int):
        self.max_input = max_input
        self.max_output = max_output

    def check(self, input_tokens: int, output_tokens: int) -> bool:
        return (input_tokens <= self.max_input and
                output_tokens <= self.max_output)

    def truncate_input(self, text: str, encoding) -> str:
        tokens = encoding.encode(text)
        if len(tokens) <= self.max_input:
            return text
        return encoding.decode(tokens[:self.max_input])

# Nutzung
budget = TokenBudget(max_input=4096, max_output=512)

if not budget.check(input_tokens, request_max):
    print(f"Input {input_tokens} ueberschreitet Budget {budget.max_input}")
    text = budget.truncate_input(text, encoding)
\end{lstlisting}

% ============================================================
% ANTI-PATTERNS
% ============================================================
\section{Anti-Patterns}

\begin{itemize}[leftmargin=*]
    \item \textbf{Kein Token-Counting}: Ohne Zählung weißt du nie, wie viel deine Anfragen kosten oder wie nah du am Context-Limit bist. Immer vor dem API-Call zählen!
    \item \textbf{max\_tokens zu groß setzen}: 4096 max\_tokens bei einer binären Klassifikation ist Verschwendung. Output-Tokens sind 3--10x teurer als Input.
    \item \textbf{Bilder als Base64 ohne Optimierung}: Detail-Level (low/high/auto) reduziert Token-Kosten für Bilder um Faktor 50. \texttt{detail: "low"} reicht für die meisten Use Cases.
    \item \textbf{History nie komprimieren}: Ein Chat mit 100 Nachrichten kostet schnell \$1 pro Request durch die kumulierte History.
    \item \textbf{Prompt-Caching nicht nutzen}: Wenn der gleiche System-Prompt 10.000x gesendet wird, aber nie gecached, zahlst du jedes Mal neu.
    \item \textbf{Auf Tokens optimieren statt auf Ergebnisse}: Ein extrem komprimierter Prompt spart Tokens, aber wenn die Ergebnisse schlechter werden, ist die Optimierung nutzlos. Immer Quality-Tests machen!
    \item \textbf{Provider-übergreifend Tokens vergleichen}: 100 Tokens bei OpenAI sind nicht gleich 100 Tokens bei Anthropic. Jeder Tokenizer zählt anders.
\end{itemize}

\begin{praxishinweis}
Das häufigste Anti-Pattern: maximale Kompression ohne Qualitätskontrolle. Ein extrem verdichteter Prompt spart Tokens, aber wenn die Ergebnisse leiden, ist die Einsparung ein Verlustgeschäft. Miss beides: Token-Kosten und Output-Qualität.
\end{praxishinweis}

% ============================================================
% PERFORMANCE
% ============================================================
\section{Performance und Optimierung}

\subsection{Latenz durch Token-Anzahl}

Die Token-Zahl beeinflusst direkt die Latenz:

\begin{table}[H]
\centering
\begin{tabularx}{\linewidth}{lXX}
\toprule
\textbf{Token-Bereich} & \textbf{Time-to-First-Token} & \textbf{Gesamtzeit (500 Out)} \\
\midrule
1K Input & 200--400ms & 1,2s \\
10K Input & 500--800ms & 1,5s \\
100K Input & 2--4s & 3--5s \\
\midrule
\end{tabularx}
\caption{Latenz in Abhängigkeit der Input-Token-Zahl (GPT-4o)}
\label{tab:latenz-tokens}
\end{table}

\subsection{Caching-Strategien}

\begin{itemize}[leftmargin=*]
    \item \textbf{Prompt-Caching (Provider-seitig)}: Automatisch bei OpenAI (ab 1K), manuell bei Anthropic. Kostet 50--90\,\% weniger für gecachte Input-Tokens.
    \item \textbf{Response-Caching (Applikations-seitig)}: Gleiche Anfragen (identischer Prompt) werden gecached. Wirksam bei Classification und Extraction.
    \item \textbf{Embedding-Caching}: Berechnete Embeddings für Dokumente cachen, statt sie neu zu berechnen.
\end{itemize}

\begin{lstlisting}[language=Python, caption=Response-Caching mit Disk-Cache]
import hashlib
import json
import os

class ResponseCache:
    def __init__(self, cache_dir: str = ".token_cache"):
        self.cache_dir = cache_dir
        os.makedirs(cache_dir, exist_ok=True)

    def _key(self, messages: list[dict], model: str) -> str:
        raw = json.dumps({"messages": messages, "model": model},
                         sort_keys=True)
        return hashlib.sha256(raw.encode()).hexdigest()

    def get(self, messages, model):
        key = self._key(messages, model)
        path = f"{self.cache_dir}/{key}.json"
        if os.path.exists(path):
            with open(path) as f:
                return json.load(f)
        return None

    def set(self, messages, model, response):
        key = self._key(messages, model)
        with open(f"{self.cache_dir}/{key}.json", "w") as f:
            json.dump(response, f)
\end{lstlisting}

\subsection{Kosten-Tracking in Produktion}

Jeder API-Call sollte geloggt werden:

\begin{lstlisting}[language=Python, caption=Token-Kosten-Tracking]
from dataclasses import dataclass
from datetime import datetime
import sqlite3

@dataclass
class TokenUsage:
    provider: str
    model: str
    input_tokens: int
    output_tokens: int
    cached_tokens: int
    cost: float
    user_id: str
    endpoint: str
    timestamp: datetime

class CostTracker:
    def __init__(self, db_path: str = "costs.db"):
        self.conn = sqlite3.connect(db_path, check_same_thread=False)
        self.conn.execute("""
            CREATE TABLE IF NOT EXISTS token_usage (
                id INTEGER PRIMARY KEY AUTOINCREMENT,
                provider TEXT, model TEXT,
                input_tokens INTEGER, output_tokens INTEGER,
                cached_tokens INTEGER DEFAULT 0,
                cost REAL, user_id TEXT, endpoint TEXT,
                timestamp TEXT
            )
        """)
        self.conn.commit()

    def log(self, usage: TokenUsage):
        self.conn.execute(
            "INSERT INTO token_usage "
            "(provider, model, input_tokens, output_tokens, "
            " cached_tokens, cost, user_id, endpoint, timestamp) "
            "VALUES (?, ?, ?, ?, ?, ?, ?, ?, ?)",
            (usage.provider, usage.model,
             usage.input_tokens, usage.output_tokens,
             usage.cached_tokens, usage.cost,
             usage.user_id, usage.endpoint,
             usage.timestamp.isoformat())
        )
        self.conn.commit()

    def daily_cost(self, date: str) -> float:
        row = self.conn.execute(
            "SELECT SUM(cost) FROM token_usage WHERE date(timestamp) = ?",
            (date,)
        ).fetchone()
        return row[0] or 0.0

    def top_users(self, limit: int = 10) -> list[tuple]:
        return self.conn.execute(
            "SELECT user_id, SUM(cost) as total "
            "FROM token_usage "
            "GROUP BY user_id ORDER BY total DESC LIMIT ?",
            (limit,)
        ).fetchall()
\end{lstlisting}

% ===========================================================%
% SICHERHEIT
% ===========================================================%
\section{Sicherheit}

\subsection{Token-Leakage vermeiden}

Token-Zahlen können Informationen über den Prompt preisgeben: Zwei sehr ähnliche Prompts erzeugen ähnliche Token-Zahlen. Bei der Logging-Architektur muss klar sein:

\begin{itemize}[leftmargin=*]
    \item \textbf{Token-Zahlen sind keine PII}, aber sie können Rückschlüsse auf Prompt-Längen und -Komplexität zulassen.
    \item \textbf{Keine Rohdaten in Logs}: Logge Token-Zahlen, aber nicht den Prompt-Text selbst, wenn dieser sensible Daten enthält.
    \item \textbf{Aggregierte Metriken}: Fasse Kosten pro User/Team/Abteilung zusammen, statt Einzel-Requests zu loggen.
\end{itemize}

\subsection{Denial-of-Wallet}

Ohne Budget-Limits kann ein einzelner User oder ein fehlerhafter Prozess das gesamte API-Budget aufbrauchen:

\begin{lstlisting}[language=Python, caption=Denial-of-Wallet-Praevention]
class BudgetGuard:
    def __init__(self, max_daily_cost: float = 100.0):
        self.max_daily_cost = max_daily_cost
        self.tracker = CostTracker()

    def check_budget(self) -> bool:
        today = datetime.now().strftime("%Y-%m-%d")
        cost = self.tracker.daily_cost(today)
        return cost < self.max_daily_cost

    def guard(self, func):
        """Decorator: Blockiert API-Calls, wenn Budget erreicht."""
        def wrapper(*args, **kwargs):
            if not self.check_budget():
                raise RuntimeError(
                    f"Daily budget {self.max_daily_cost}$ reached"
                )
            return func(*args, **kwargs)
        return wrapper
\end{lstlisting}

\subsection{Multi-Tenant-Token-Isolation}

In SaaS-Systemen müssen Token-Kosten pro Mandant (Tenant) getrennt werden:

\begin{itemize}[leftmargin=*]
    \item \textbf{Separate API-Keys pro Tenant}: Ermöglicht individuelle Limits und Abrechnung
    \item \textbf{Per-Tenant-Budget}: Jeder Tenant hat sein eigenes max\_daily\_cost-Limit
    \item \textbf{Cost-Allocation}: Token-Verbrauch wird auf Tenant-Kostenstellen gebucht (chargeback-fähig)
\end{itemize}

% ===========================================================
% ZUSAMMENFASSUNG (existiert, erweitert)
% ===========================================================
\section{Zusammenfassung}

\begin{itemize}[leftmargin=*]
    \item Ein Token ist kein Wort -- deutsche Texte produzieren etwa 1,4--2x mehr Tokens als englische (Komposita!)
    \item TikToken ist das Werkzeug zur Token-Zählung: Ohne es tappst du im Dunkeln bei Kosten und Limits
    \item max\_tokens kontrolliert nicht nur die Antwortlänge, sondern ist ein direkter Kostenhebel -- setze ihn bewusst!
    \item Prompt-Kompaktierung und Chat-History-Management sind die effektivsten Token-Einsparer
    \item Prompt-Caching (OpenAI + Anthropic) kann Input-Kosten um bis zu 90\,\% reduzieren
    \item Output-Tokens sind 3--10x teurer als Input-Tokens -- priorisiere kurze Antworten
    \item Token-Budgeting und Cost-Tracking sind Grundvoraussetzungen für produktive Systeme
\end{itemize}

% ===========================================================
% MERKE-BOX
% ===========================================================
\begin{merke}
\begin{itemize}
    \item \textbf{Kosten liegen im Output}: Output-Tokens sind 3--10x teurer als Input. max\_tokens ist der wichtigste Kostenhebel.
    \item \textbf{Deutsch ist teurer}: Deutsche Texte kosten 30--50\,\% mehr als englische durch längere Tokenisierung.
    \item \textbf{Caching ist Pflicht}: Prompt-Caching spart 50--90\,\% der Input-Kosten. Immer aktivieren!
    \item \textbf{History killt Budget}: Ohne Kompression wächst die Chat-History unbegrenzt -- komprimiere ab 80\,\% des Context-Limits.
    \item \textbf{Token-Zahlen sind nicht portabel}: 100 Tokens bei OpenAI $\neq$ 100 Tokens bei Anthropic.
\end{itemize}
\end{merke}

% ===========================================================
% PRAXISPROJEKT
% ===========================================================
\section{Praxisprojekt}

\subsection{Aufgabe: Multi-Provider Cost Dashboard}

Baue ein System, das die Token-Kosten über mehrere Provider hinweg trackt und visualisiert.

\textbf{Anforderungen:}
\begin{itemize}[leftmargin=*]
    \item Implementiere einen \texttt{TokenTracker}, der API-Calls über OpenAI, Anthropic und Google loggt
    \item Verwende die \texttt{CostTracker}-Klasse aus diesem Kapitel als Basis
    \item Füge einen \texttt{BudgetGuard} hinzu, der bei Überschreitung des Tagesbudgets automatisch auf ein günstigeres Modell fallbackt ("GPT-4o $\rightarrow$ GPT-4o-mini")
    \item Implementiere ein Reporting, das die Top-5-Token-Verbraucher pro Tag ausgibt
    \item Schreibe einen Test, der die Kosten bei 1.000 Requests mit und ohne Prompt-Caching vergleicht
\end{itemize}

\textbf{Testdaten:} Simuliere 1.000 API-Calls mit variierenden Input-Längen (50--4.000 Tokens) und Output-Längen (10--500 Tokens).

\textbf{Erfolgskriterium:} Das Dashboard zeigt korrekte Tageskosten, warnt bei Budget-Überschreitung und falls automatisch auf das günstigere Modell zurück.

% ===========================================================
% WEITERFÜHRENDE RESSOURCEN
% ===========================================================
\section{Weiterführende Ressourcen}

\subsection{Dokumentationen}

\begin{itemize}[leftmargin=*]
    \item OpenAI Tokenizer: \href{https://platform.openai.com/tokenizer}{platform.openai.com/tokenizer}
    \item TikToken GitHub: \href{https://github.com/openai/tiktoken}{github.com/openai/tiktoken}
    \item Anthropic Tokenizer: \href{https://docs.anthropic.com/en/docs/build-with-claude/token-counting}{docs.anthropic.com}
    \item Google Token Counting: \href{https://ai.google.dev/gemini-api/docs/tokens}{ai.google.dev/gemini-api/docs/tokens}
\end{itemize}

\subsection{Tools}

\begin{itemize}[leftmargin=*]
    \item \textbf{TikToken}: \href{https://github.com/openai/tiktoken}{github.com/openai/tiktoken} - BPE-Tokenizer für OpenAI-Modelle
    \item \textbf{anthropic-tokenizer}: \href{https://github.com/anthropics/anthropic-tokenizer}{github.com/anthropics/anthropic-tokenizer}
    \item \textbf{Langfuse}: \href{https://langfuse.com}{langfuse.com} - Open-Source-Tracking für LLM-Kosten
    \item \textbf{LiteLLM}: \href{https://github.com/BerriAI/litellm}{github.com/BerriAI/litellm} - Einheitliche Kostenberechnung über alle Provider
\end{itemize}

\subsection{Bücher und Artikel}

\begin{itemize}[leftmargin=*]
    \item OpenAI Cookbook: \href{https://cookbook.openai.com}{cookbook.openai.com} - "How to count tokens with tiktoken"
    \item Sennrich et al. (2016): "Neural Machine Translation of Rare Words with Subword Units" -- grundlegende BPE-Arbeit
    \item Kudo \& Richardson (2018): "SentencePiece: A simple and language independent subword tokenizer"
\end{itemize}

\textbf{Nächster Schritt:} Mit diesem Token-Wissen im Kopf geht es weiter zum Kernstück der AI Engineer-Arbeit: Prompt Engineering.
